%% ============================================================================
%% Abstract and keywords (independent file, included inside \begin{frontmatter})
%% ============================================================================

\begin{abstract}
Sparse-view cone-beam CT (CBCT) reconstruction aims to recover a tomographic
volume from an extremely limited set of projection views, which is an
ill-posed inverse problem that traditional methods make it difficult to 
achieve effective reconstruction, and is typically solved with deterministic
deep networks. A probabilistic formulation can in principle quantify
reconstruction uncertainty and exploit privileged information --- dense views
or even the ground-truth volume --- during training, as long as inference
never depends on it. In this work we present the \emph{GT-conditioned CVAE}
 a conditional variational extension of the 3D-Gaussian-splatting CBCT
reconstructor. During training only, a lightweight 3D convolutional network
encodes the ground-truth volume into a 32-dimensional descriptor that
conditions the variational posterior $q(z\mid C,Y)$ over a 64-dimensional
global scene latent, while a sparse-view neural-process prior
$p(z\mid C)$ is trained jointly. At deployment the GT encoder and the
posterior are discarded: $z$ is drawn exclusively from the sparse-view prior,
so inference requires exactly the same inputs as the deterministic backbone
and additionally yields per-point predictive variance from $N{=}8$
Monte-Carlo samples. The decoder renders exclusively from the sparse
context features, so the latent is the sole channel of dense information
during training. On the LUNA16 CBCT benchmark with six
projections, the GT-conditioned CVAE reaches $29.50$\,dB PSNR /
$0.868$\,SSIM, within $0.04$\,dB of the deterministic backbone
($29.54$\,dB / $0.870$\,SSIM), while providing calibrated deployment-side
uncertainty estimates at no additional inference cost. Ablations
quantify the decisive design choices -- posterior access to the ground
truth and fusion token initialization -- for future probabilistic CBCT
models.
\end{abstract}

%% Research highlights
\begin{highlights}
\item A GT-conditioned CVAE for Gaussian-splatting sparse-view CBCT
reconstruction: ground truth is encoded into the variational posterior
during training only.
\item Deployment discards both the GT encoder and dense views, sampling
exclusively from a sparse-view prior; per-point predictive variance comes
for free from Monte-Carlo estimation.
\item Ablations isolate the decisive design choices: posterior access to
the ground truth and fusion token initialization ($\sim\!2$\,dB), with the
decoder routed through the context features.
\end{highlights}

%% Keywords
\begin{keyword}
Sparse-view CBCT reconstruction \sep 3D Gaussian splatting \sep
conditional variational autoencoder \sep variational inference \sep
latent neural processes \sep uncertainty estimation
\end{keyword}